\section{Risultati}
\label{sec:results}

Sono stati eseguiti 16 esperimenti su una griglia comune di configurazioni
(numero di client $N \in \{3,5,8\}$, fanout $k$, passi locali $H$),
metà in esecuzione locale e metà su
AWS EC2 in modalità multi-istanza.
In locale è stata sfruttata una gpu nvidia, su aws no. Per ogni
configurazione sono riportati: l'accuratezza di validazione media tra i
worker, l'accuratezza sul global test set (calcolata al round del checkpoint
migliore su validation di ciascun worker) con il
relativo \emph{spread} (differenza massimo$-$minimo tra worker, in punti
percentuali), la distanza L2 media tra i pesi del best model dei worker, il numero
totale di messaggi scambiati e il wall-clock di sistema. La
Tabella~\ref{tab:master} riporta tutti e 16 i run.

Questa griglia non è un fattoriale completo su $N$, $k$ e $H$: abbiamo adottato un apporccio a blocchi per ridurre il numero di configurazioni.
È stato quindi adottato un disegno sequenziale a
blocchi, un fattore alla volta: il Blocco A (3 run, $N=3$, $k=1$) isola
l'effetto di $H \in \{100,500,1000\}$, individuando $H=1000$ come migliore per ognuna delle metriche;
il Blocco B (1 run) isola un piccolo aumento di $k$; il Blocco C (4 run,
$H=1000$) incrocia $N \in \{5,8\}$ con $k \in \{1, N{-}1\}$, i due estremi
del fanout, isolando l'effetto congiunto di scala e grado di connettività.

\begin{table*}[t]
  \centering
  \caption{Risultati per tutte le 16 configurazioni sperimentali ($k$=fanout, $H$=passi locali; spread e L2 sul global test set — Sez.~\ref{sec:results}).}
  \label{tab:master}
  \scriptsize
  \begin{tabular}{llrrrrrrrrr}
    \toprule
    \textbf{Config} & \textbf{Ambiente} & $N$ & $k$ & $H$ & \textbf{Val.\ acc.} & \textbf{Global-test acc.} & \textbf{Spread} & \textbf{L2} & \textbf{Msg inviati} & \textbf{Wall-clock} \\
    \midrule
    ref      & locale & 3 & 1 & 500  & 0.898$\pm$0.010 & 0.769$\pm$0.101 & 19.1\% & 102.9 & 88   & 5.2 min \\
    h100     & locale & 3 & 1 & 100  & 0.874$\pm$0.024 & 0.760$\pm$0.126 & 24.9\% & 235.4 & 66   & 5.0 min \\
    h1000    & locale & 3 & 1 & 1000 & 0.906$\pm$0.007 & 0.786$\pm$0.091 & 17.9\% & 184.3 & 87   & 8.7 min \\
    f2       & locale & 3 & 2 & 1000 & 0.890$\pm$0.017 & 0.784$\pm$0.089 & 15.4\% & 107.7 & 172  & 8.1 min \\
    n5\_f1   & locale & 5 & 1 & 1000 & 0.914$\pm$0.011 & 0.773$\pm$0.110 & 21.2\% & 90.7  & 228  & 15.8 min \\
    n5\_fn   & locale & 5 & 4 & 1000 & 0.908$\pm$0.025 & 0.792$\pm$0.104 & 21.8\% & 117.0 & 491  & 12.6 min \\
    n8\_f1   & locale & 8 & 1 & 1000 & 0.910$\pm$0.016 & 0.760$\pm$0.098 & 22.4\% & 174.0 & 295  & 19.9 min \\
    n8\_fn   & locale & 8 & 7 & 1000 & 0.905$\pm$0.032 & 0.835$\pm$0.068 & 15.2\% & 94.7  & 1863 & 24.0 min \\
    \addlinespace
    n3\_h500  & AWS & 3 & 1 & 500  & 0.901$\pm$0.006 & 0.768$\pm$0.086 & 16.8\% & 105.8 & 99   & 138.9 min \\
    n3\_h100  & AWS & 3 & 1 & 100  & 0.886$\pm$0.007 & 0.738$\pm$0.081 & 14.5\% & 74.9  & 165  & 158.3 min \\
    n3\_h1000 & AWS & 3 & 1 & 1000 & 0.905$\pm$0.015 & 0.782$\pm$0.099 & 18.8\% & 91.5  & 96   & 135.2 min \\
    n3\_f2    & AWS & 3 & 2 & 1000 & 0.887$\pm$0.016 & 0.784$\pm$0.065 & 12.0\% & 66.1  & 184  & 120.4 min \\
    n5\_f1    & AWS & 5 & 1 & 1000 & 0.915$\pm$0.019 & 0.773$\pm$0.119 & 22.5\% & 101.0 & 232  & 223.9 min \\
    n5\_fn    & AWS & 5 & 4 & 1000 & 0.902$\pm$0.015 & 0.785$\pm$0.079 & 19.5\% & 98.0  & 572  & 200.6 min \\
    n8\_f1    & AWS & 8 & 1 & 1000 & 0.910$\pm$0.025 & 0.754$\pm$0.122 & 30.1\% & 223.1 & 320  & 207.8 min \\
    n8\_fn    & AWS & 8 & 7 & 1000 & 0.896$\pm$0.027 & 0.816$\pm$0.061 & 15.8\% & 64.7  & 1303 & 185.7 min \\
    \bottomrule
  \end{tabular}
\end{table*}

\subsection{Scalabilità: accuratezza e tempo vs.\ $N$}
La Fig.~\ref{fig:accuracy_vs_n} mostra accuratezza di validazione e
global-test al variare di $N$, per $H=1000$ fissato, confrontando $k=1$
con $k=N-1$. L'accuratezza di
validazione resta stabile tra l'89\% e il 92\% per ogni $N$ testato, con
variazioni contenute entro circa due punti percentuali tra $k=1$ e $k=N-1$: il
sistema scala senza perdita di qualità sul dato visto localmente. Sul
global test set l'effetto di $k$ è più marcato: a $N=8$, passare da $k=1$ a
$k=N-1$ porta l'accuratezza global-test dal 76,0\% all'83,5\% in locale
e dal 75,4\% all'81,6\% su AWS — un guadagno di
6--8 punti percentuali attribuibile a un numero maggiore di aggiornamenti
ricevuti per round.

\begin{figure}[t]
  \centering
  \includegraphics[width=\linewidth]{accuracy_vs_n}
  \caption{Accuratezza di validazione (sinistra) e sul global test set
  (destra) al variare di $N$, per $k=1$ e $k=N-1$, locale vs.\ AWS.}
  \label{fig:accuracy_vs_n}
\end{figure}

\subsection{Overhead di comunicazione}
Ogni messaggio push trasporta l'intero \texttt{state\_dict} del modello. Il traffico totale
stimato (Fig.~\ref{fig:communication}) cresce con $N$ e, soprattutto, con
$k$: a $N=8$, passare da $k=1$ a $k=N-1$ porta il numero di messaggi
scambiati da 295 a 1863 in locale e da 320 a 1303 su AWS. Questo è il costo diretto, e prevedibile, della scelta di $k$:
l'aumento di accuratezza global-test osservato con $k$ alto
(Fig.~\ref{fig:accuracy_vs_n}) va quindi letto come un compromesso esplicito
comunicazione/qualità.

\begin{figure}[t]
  \centering
  \includegraphics[width=\linewidth]{communication_overhead}
  \caption{Traffico totale stimato (messaggi $\times$ 6,82\,MB) per
  configurazione.}
  \label{fig:communication}
\end{figure}

\subsection{Dove viene speso il tempo}
La Fig.~\ref{fig:phase_timing} scompone la durata media di un round nelle tre
fasi definite in Sezione~\ref{sec:details}B. In locale la Fase A  dura 5--13\,s, la Fase B  0,3--6,9\,s in
funzione di $H$, la Fase C resta sotto i 340\,ms anche con
$k=N-1$. Su AWS la Fase A si stabilizza attorno a 109--126\,s
indipendentemente da $N$ e da $H$ — un valore pressoché costante compatibile
con un collo di bottiglia sul throughput di inferenza delle istanze
\texttt{t3.small} piuttosto che con un effetto del
workload. La Fase C, invece, resta nello stesso ordine di grandezza del caso
locale (60--227\,ms) anche attraverso rete reale tra istanze EC2 separate:
l'overhead di comunicazione del protocollo di push non è la causa del
rallentamento osservato su AWS.

\begin{figure}[t]
  \centering
  \includegraphics[width=\linewidth]{phase_timing}
  \caption{Durata media delle fasi A, B, C per round, locale vs.\ AWS.}
  \label{fig:phase_timing}
\end{figure}

\subsection{Locale vs.\ AWS multi-istanza}
Per le 8 configurazioni condivise dai due ambienti, la
Fig.~\ref{fig:local_vs_aws} confronta wall-clock di sistema e spread
global-test. Il wall-clock su AWS è 8--32$\times$ superiore a quello locale
per la stessa configurazione (dominato dalla Fase A, come discusso sopra).
Lo spread global-test resta invece in una banda solo in parte sovrapposta tra
i due ambienti: 12--30\% su AWS contro 15--25\% in locale — un intervallo
più ampio lato AWS, trainato dal 30,1\% di \texttt{n8\_f1} (il valore più
alto misurato in tutta la campagna) e dal 12,0\% di \texttt{n3\_f2} (il più
basso). Non emerge quindi un vantaggio sistematico di un ambiente
sull'altro sulla qualità della convergenza funzionale, ma la dispersione tra
configurazioni è più marcata su AWS in questa campagna.

\begin{figure}[t]
  \centering
  \includegraphics[width=\linewidth]{local_vs_aws}
  \caption{Wall-clock di sistema (scala log) e spread di accuratezza
  global-test, per le 8 configurazioni condivise da locale e AWS.}
  \label{fig:local_vs_aws}
\end{figure}

\subsection{Analisi per classe: matrice di confusione (campagna AWS)}
La campagna AWS qui riportata calcola, oltre ad accuratezza aggregata e
spread, anche precision/recall/F1 macro (media semplice tra le 62 classi,
non pesata per il numero di campioni di ciascuna — a differenza
dell'accuratezza aggregata) e la matrice di confusione completa
sul global test set per ciascuna configurazione. La
Tabella~\ref{tab:macro} riporta le metriche macro medie tra worker, con lo
stesso criterio [B] (checkpoint a minima validation loss) usato per
l'accuratezza in Tabella~\ref{tab:master}.

\begin{table*}[t]
  \centering
  \caption{Metriche macro sul global test set, media tra worker, per le 8 configurazioni AWS (criterio [B]).}
  \label{tab:macro}
  \scriptsize
  \begin{tabular}{lrrrrrr}
    \toprule
    \textbf{Config} & $N$ & $k$ & \textbf{Macro-precision} & \textbf{Macro-recall} & \textbf{Macro-F1} & \textbf{Spread F1} \\
    \midrule
    n3\_h500  & 3 & 1 & 0.750 & 0.710 & 0.685 & 11.2\% \\
    n3\_h100  & 3 & 1 & 0.719 & 0.678 & 0.643 & 7.2\%  \\
    n3\_h1000 & 3 & 1 & 0.764 & 0.730 & 0.698 & 13.7\% \\
    n3\_f2    & 3 & 2 & 0.757 & 0.747 & 0.715 & 10.0\% \\
    n5\_f1    & 5 & 1 & 0.772 & 0.732 & 0.704 & 20.1\% \\
    n5\_fn    & 5 & 4 & 0.775 & 0.742 & 0.714 & 15.5\% \\
    n8\_f1    & 8 & 1 & 0.758 & 0.708 & 0.679 & 28.6\% \\
    n8\_fn    & 8 & 7 & 0.786 & 0.752 & 0.742 & 13.2\% \\
    \bottomrule
  \end{tabular}
\end{table*}

Il macro-F1 segue lo stesso ordinamento qualitativo dell'accuratezza
aggregata (\texttt{n8\_fn} migliore, \texttt{n8\_f1} e \texttt{n3\_h100}
peggiori) ma è sistematicamente più basso di essa (0,64--0,74 contro
0,74--0,82 di accuratezza): la differenza è attesa, perché l'accuratezza
aggregata pesa ogni campione allo stesso modo mentre il macro-F1 pesa ogni
\emph{classe} allo stesso modo — su FEMNIST, dove alcune classi hanno
migliaia di campioni e altre poche centinaia (Sezione~\ref{sec:background}C),
un'accuratezza aggregata alta può convivere con prestazioni scarse sulle
classi rare, ed è esattamente ciò che il macro-F1 rivela.

Sommando le matrici di confusione dei worker di ciascuna run, lo schema degli
errori è pressoché identico nelle 8 configurazioni: le stesse tre coppie di
caratteri dominano gli errori in tutti e 8 i run — \texttt{l}$\to$\texttt{1}
(56--74\% delle occorrenze della vera classe \texttt{l}),
\texttt{O}$\to$\texttt{0} (52--73\%) e \texttt{I}$\to$\texttt{1}
(44--62\%) — e le classi con recall più basso sono, senza eccezioni tra le
8 run, \texttt{o} (lettera minuscola), \texttt{l} e \texttt{q}; \texttt{I}
è tra le peggiori in 7 delle 8 configurazioni (l'eccezione è
\texttt{n8\_fn}, dove al suo posto compare \texttt{p}).
Il caso più estremo è la lettera minuscola \texttt{o}: recall tra lo 0\% e
lo 0,9\% in tutte le configurazioni, con il 94\% circa dei suoi campioni
classificati come cifra \texttt{0} (dato dalla run \texttt{n8\_fn}, la più
accurata della campagna). In pratica il modello collassa la terna
\texttt{0}/\texttt{O}/\texttt{o} quasi interamente sulla cifra e la terna
\texttt{1}/\texttt{l}/\texttt{I} quasi interamente sulla cifra \texttt{1} —
una conferma empirica diretta dell'ambiguità visiva già anticipata come
motivazione del label smoothing (Sezione~\ref{sec:details}E): sono
esattamente le stesse coppie di caratteri citate lì come esempio. La
consistenza dello schema attraverso tutte le configurazioni suggerisce che
non sia un effetto della topologia di comunicazione o della scala $N$, ma
un limite del compito stesso (ambiguità geometrica intrinseca tra alcune
classi di FEMNIST) che nessuna delle configurazioni testate risolve.
